\documentclass[10pt]{article}
\usepackage[margin=2.2cm]{geometry}
\usepackage{graphicx}
\usepackage{amsmath}
\usepackage{booktabs}
\usepackage[hidelinks]{hyperref}
\setlength{\parskip}{4pt}\setlength{\parindent}{0pt}

\title{\vspace{-1.2cm}Neuron mechanisms III: where the degree-16 parity lives\\
\large \texttt{wd3\_k0\_N2048}, step 500k --- post-nonlinearity units, uniform evaluation data}
\author{}\date{September 2026}

\begin{document}\maketitle\vspace{-1.4cm}

\section*{Setup}
We use the same protocol as in the $d4$/$d8$ reports: for each neuron,
$a=\phi(z)$ with $\phi=\mathrm{relu}(\tanh)$ and $z=(Vh)_n$, evaluated on all
$2^{16}$ uniform inputs. A neuron is called \emph{$d16$-dominant} when
$\chi_{d16}$ is the largest among its 15 tracked tree-mode Walsh coefficients.
To understand what drives each neuron, we decompose $\mathrm{Var}(z)$ into
contributions from the 16 raw singleton inputs, the tree degree-2/4 modes, the two
half-parities $y_0=\chi_{d8_0},\,y_1=\chi_{d8_1}$, and $\chi_{d16}$ itself.

\section*{Headline: the target parity has no dedicated neurons}
Only \textbf{25 of 8192} neurons are $d16$-dominant (6/10/3/6 per block), but
these neurons do not carry an appreciable fraction of the actual $d16$
computation.

\textbf{(i) 19 wide raw-bit gates (blocks 0--2).} Their pre-activation $z$ is built
almost entirely from \emph{all sixteen} singleton inputs ($\geq99\%$ of
$\mathrm{Var}\,z$). They are the width-16 version of the same balanced-input
majority-gate mechanism seen at lower degrees. At width 16, however, the neuron's
Walsh spectrum is distributed over so many modes that each individual tree
coefficient is extremely small. Their tree coverage is only $0.007$, continuing
the sequence
$d2\!:\!1/3 \to d4\!:\!0.15 \to d8\!:\!0.044 \to d16\!:\!0.007$, and their
$d16$ power is only $a^2(d16)\sim10^{-5}$--$10^{-4}$. Thus their apparent
``$d16$-dominance'' mostly comes from taking an argmax over fifteen coefficients
that are all close to zero.

\textbf{(ii) 6 near-dead receivers (block 3).} These neurons do have a genuine
$d16$-shaped component: about $40\%$ of their activation variance is in
$\chi_{d16}$, and their tree coverage is $\approx0.95$. But their total activation
variance is only $\mathrm{Var}(a)\sim10^{-6}$, so their contribution to the
network-level computation is negligible.

\section*{Who actually carries $\chi_{d16}$}
The layer-level accounting over \emph{all} neurons (Fig.~\ref{fig:d16}, middle)
makes the answer clear. The total degree-16 power $H(d16)$ is
$1.6{\times}10^{-3}$ and $1.4{\times}10^{-3}$ in blocks 0--1, which is essentially
the noise floor. It then rises to $2.8{\times}10^{-2}$ in block 2 and
$1.3{\times}10^{-1}$ in block 3.

At exactly these blocks, almost all of the $d16$ signal is carried by the
\textbf{$8{+}8$ corner units identified in the $d8$ study}: they account for
$98.7\%$ of $H(d16)$ in block 2 and $99.8\%$ in block 3. There are 189 such
neurons. Their input has the form
$\gamma(\pm y_0\pm y_1)$, where $y_0$ and $y_1$ are the two degree-8 half-parities,
so after the nonlinearity they act as two-input corner detectors. As in the
degree-2 case, each such corner detector has equal nonconstant Walsh power in its
two inputs and in their product. Here that means equal thirds over
$\{\chi_{d8_0},\chi_{d8_1},\chi_{d16}\}$.

By contrast, the 25 neurons labeled ``$d16$-dominant'' contribute only $0.2\%$ of
$H(d16)$ in block 2 and $0.04\%$ in block 3.

Unlike $d8$ --- which is first created in blocks 0--1 and then passed to dedicated
\emph{amplifier} neurons --- no analogous $d16$ amplifier population appears.
The reason is architectural: $d16$ is created only in blocks 2--3, immediately
before the final linear readout uses the residual stream. There is therefore no
later residual block that needs to extract and amplify a cleaner copy of the
feature.

As a result, the degree-16 parity remains a \emph{collective} representation all
the way to the readout. It appears as one component of many $8{+}8$ corner
neurons, rather than becoming the main computation of a dedicated neuron.

\begin{figure}[h]\centering
\includegraphics[width=\textwidth]{fig5_d16_mechanisms.pdf}
\caption{Left: the 25 $d16$-dominant neurons in input space --- 19 wide raw-bit gates
at the origin, 6 near-dead receivers; the actual creators (red, $d8$-labeled corner
units) sit at $x\!\approx\!1$. Middle: share of layer $H(d16)$ by class --- at the
blocks where $H(d16)$ is macroscopic ($2.8{\times}10^{-2}$, $1.3{\times}10^{-1}$), the
corner units own it ($98.7\%$, $99.8\%$). Right: per-neuron $a^2(d16)$ by class ---
the strongest corners reach $10^{-2}$, two orders above anything
``$d16$-dominant''.}
\label{fig:d16}
\end{figure}

\begin{table}[h]\centering\small
\begin{tabular}{lccccc}
\toprule
class & input ($z$) & blocks & $n$ & tree cov. & share of $H(d16)$ @ b2/b3 \\
\midrule
$8{+}8$ corner ($d8$-labeled) & $\gamma(\pm y_0\pm y_1)$ & 2--3 & 189 & $1.00$ & $98.7\%$ / $99.8\%$ \\
16-input raw-bit gate & $\gamma\sum_{16}s_kx_k$ & 0--2 & 19 & $0.007$ & $0.2\%$ / --- \\
near-dead receiver & mixed $d16{+}d8$, tiny & 3 & 6 & $0.95$ & --- / $0.04\%$ \\
\bottomrule
\end{tabular}
\end{table}

\textbf{Takeaway.} The mechanism across the full parity hierarchy is now
consistent. Degree-2, 4, and 8 modes are generated by the same basic template:
balanced linear combinations passed through $\mathrm{relu}(\tanh)$. When the
inputs to such a gate are raw bits, its Walsh power is spread over many modes, so
the fraction visible in the target-tree modes decreases rapidly with width
($1/3,\,0.15,\,0.044,\,0.007$).

The only genuinely compositional step occurs at the top of the hierarchy. The two
degree-8 half-parities are already available, and $8{+}8$ corner detectors combine
them to produce $\chi_{d16}$. In each such neuron, the degree-16 product is one of
three equal nonconstant Walsh components rather than a dedicated feature. Because
the readout follows immediately afterward, the network never needs to distill
$\chi_{d16}$ into a separate amplifier population.

Therefore, in the meaningful mechanistic sense, ``the neurons computing the
$d16$ parity'' are the $8{+}8$ corner-detector population as a whole. The neurons
selected purely by a $d16$-dominance label mostly reflect an argmax over tiny
coefficients rather than the true location of the computation.

\vfill\footnotesize
Artifacts: \texttt{gdrive:parity/analysis\_output/neuron\_mechanisms/wd3\_k0\_N2048/step\_00500000/}
(\texttt{d16\_classes.csv}, figures, this report; companion $d2$/$d4$/$d8$ probes in
the same folder).
\end{document}
